\section{Conclusiones}

El presente proyecto nació de una inquietud genuina: ¿por qué, en plena era de
la información, el consumidor sigue tomando decisiones de compra en la
oscuridad? A lo largo de este semestre —a través de revisión exhaustiva de
literatura académica y la aplicación de una encuesta a una muestra
considerablemente amplia — fue tomando forma una respuesta cada vez más clara:
el problema no es la falta de información, sino su inaccesibilidad deliberada. 

La obsolescencia programada, lejos de ser un fenómeno tecnológico inevitable, es
una decisión de diseño. Como quedó documentado a lo largo del marco teórico,
aproximadamente el 80\% de la contaminación asociada a los dispositivos
electrónicos se origina en la etapa de diseño del producto —el mismo momento en
que se determina su vida útil limitada. Esto revela que el desperdicio no es una
consecuencia accidental del progreso tecnológico, sino una condición
deliberadamente integrada en él. La práctica de la incompatibilidad de
actualizaciones y la restricción de piezas de repuesto son expresiones concretas
de un sistema que prioriza el reemplazo sobre la continuidad. 

Frente a este panorama, los resultados obtenidos mediante la encuesta aplicada a
241 habitantes del municipio de Culiacán —representativos del segmento
poblacional de 27 años o más— ofrecen una lectura valiosa del estado actual del
consumidor mexicano. Los datos recopilados permitieron confirmar que existe una
brecha significativa entre la intención de consumir responsablemente y la
capacidad real de hacerlo: si bien una parte importante de los encuestados
manifestó disposición a reparar sus dispositivos bajo condiciones accesibles y
justas, la mayoría reconoció enfrentar dificultades para encontrar información
clara y confiable al momento de realizar una compra. Este hallazgo valida
directamente la hipótesis central del proyecto: el consumidor no carece de
voluntad, carece de herramientas. 

Asimismo, los resultados evidenciaron que la sobrecarga informativa —lejos de
empoderar— paraliza y aunque el consumidor no sea consciente del llamado
Greenwashing, éste es condicionado por este último sin saberlo a moldear sus
hábitos de consumo. Ambos fenómenos operan en beneficio del ciclo de consumo
acelerado que sostiene la obsolescencia programada, perpetuando una asimetría de
información que favorece sistemáticamente a los fabricantes por encima del
usuario. La sección de reacción ante estímulos publicitarios de la encuesta fue
especialmente reveladora en este sentido: los participantes demostraron ser
capaces de identificar la diferencia entre un anuncio que informa y uno que
seduce, aunque en la práctica cotidiana esa distinción raramente se traduce en
una decisión de compra más consciente —precisamente porque no existe una
herramienta que facilite ese puente. 

La respuesta que esta investigación propone ante dicho vacío es una aplicación
móvil fundamentada en el principio de divulgación progresiva, diseñada para
reducir la carga cognitiva del usuario y entregarle, a través del escaneo de un
código de barras, información centralizada, verificable y comprensible sobre el
dispositivo que desea adquirir. Esta propuesta no pretende ser únicamente una
solución tecnológica: es, en su núcleo, una herramienta de equidad informativa.
Al integrar datos de sostenibilidad, reseñas no patrocinadas, certificaciones
verificables y un directorio de centros de reparación, la aplicación busca
traducir los principios de la economía circular —reparar, restaurar, replantear—
en acciones concretas y accesibles para cualquier persona, independientemente de
su nivel de conocimiento técnico. 

En definitiva, este proyecto concluye que combatir la obsolescencia programada
no requiere únicamente de regulación legislativa ni de grandes transformaciones
industriales —aunque ambas sean necesarias. Requiere, antes que nada, de un
consumidor informado. Y un consumidor informado necesita herramientas diseñadas
honestamente para él. La aplicación propuesta representa un primer paso en esa
dirección: devolver al consumidor la autonomía que el mercado, sistemáticamente,
le ha negado. 

